% ============================================================================
% Fichier  : partie5/P05_C18_sans_fil.tex
% Titre    : Sans-Fil et Objets Connectés
% Auteur   : MINKA MI NGUIDJOÏ Thierry Emmanuel
% Version  : 1.0 -- novembre 2026
% ============================================================================
% Sources : documentation publique Wi-Fi (WPA2/WPA3) ; MITRE ATT&CK ;
% bonnes pratiques objets connectés. Entrepôt de Douala de LiZbiZ : cas
% fictif.
% ============================================================================

\chapter{Sans-Fil et Objets Connectés}
\label{chap:sansfil}
\chaptermark{Sans-fil et objets connectés}

\epigraph{%
    « Le réseau filaire s'arrête aux murs~;\\
    le réseau sans-fil franchit la rue. »%
}{Principe du cours}

\niveauChapitre{Intermédiaire}{Chapitre~\ref{chap:reseau}~; notions de
réseaux radio.}

\begin{boxobjectifs}
\begin{itemize}
    \item Décrire les faiblesses des réseaux sans-fil d'entreprise.
    \item Comprendre le point d'accès pirate et le portail jumeau.
    \item Reconnaître les risques des objets connectés.
    \item Tester un réseau sans-fil dans un cadre légal et autorisé.
    \item Rédiger les constats et les mesures de cloisonnement.
\end{itemize}
\end{boxobjectifs}

\bigskip

% ============================================================================
\section*{Question directrice : jusqu'où porte le réseau de l'entrepôt, et qui peut l'écouter~?}
% ============================================================================

L'entrepôt de Douala utilise un réseau sans-fil pour ses lecteurs de codes
et ses imprimantes d'étiquettes. Depuis le parking, Mme~NGO~MBOCK capte le
signal. Ce chapitre explique les risques du sans-fil et des objets
connectés, et comment les tester sans intercepter les communications de
tiers.

\medskip
\textbf{Plan du chapitre.}
\begin{enumerate}
    \item Les faiblesses du sans-fil (section~\ref{sec:c18-faiblesses}).
    \item Point d'accès pirate et portail jumeau (section~\ref{sec:c18-jumeau}).
    \item Les objets connectés (section~\ref{sec:c18-objets}).
    \item Tester légalement (section~\ref{sec:c18-legal}).
\end{enumerate}

\separateur

% ============================================================================
\section{Les Faiblesses du Sans-Fil}
\label{sec:c18-faiblesses}
% ============================================================================

\begin{tableaucours}{L{4.4cm}L{7.8cm}}{%
    \tableauHeader{Faiblesse} &
    \tableauHeader{Risque}
}
Clé partagée faible & Déchiffrement du trafic après capture \\
Protocole ancien & Attaques connues sur le chiffrement \\
Réseau invité relié au réseau interne & Un visiteur atteint les serveurs \\
Absence de cloisonnement & Un lecteur compromis atteint la bureautique \\
Identifiants par défaut sur les bornes & Prise de contrôle de la borne \\
\end{tableaucours}
\vspace{-0.8em}
\begin{center}{\small\itshape Tableau~18.1. Faiblesses des réseaux sans-fil}\end{center}

% ============================================================================
\section{Point d'Accès Pirate et Portail Jumeau}
\label{sec:c18-jumeau}
% ============================================================================

\begin{boxdef}
\termecle{Point d'accès pirate}~: borne non autorisée branchée sur le
réseau de l'entreprise, qui ouvre une porte dérobée.
\termecle{Portail jumeau}~: fausse borne qui imite le nom du réseau légitime
pour capter les connexions et les identifiants des usagers.
\end{boxdef}

\begin{boxalert}
\textbf{Violation typique.} Monter un portail jumeau sur le lieu de travail
et capter les connexions des salariés.

\smallskip
\textbf{Règle~:} capter les communications d'autrui est une infraction
pénale (Chapitre~\ref{chap:droit}). Un test de portail jumeau n'est permis
que sur un réseau et des comptes dédiés au test, avec accord écrit, sans
toucher aux usagers réels.
\end{boxalert}

% ============================================================================
\section{Les Objets Connectés}
\label{sec:c18-objets}
% ============================================================================

Les objets connectés (lecteurs, imprimantes, caméras, capteurs) partagent
des faiblesses récurrentes~: identifiants par défaut, micrologiciel non mis
à jour, services d'administration ouverts, absence de chiffrement. Branchés
sur le même réseau que la bureautique, ils deviennent un point d'entrée.

\begin{boxscenario}{Opération LOGONE. L'imprimante d'étiquettes}
\textbf{Constatation.} À Douala, les imprimantes d'étiquettes conservent le
mot de passe d'administration par défaut et sont sur le même réseau que les
postes. L'une d'elles expose une interface web d'administration.

\medskip
\textbf{Analyse} (confiance moyenne). Une imprimante compromise sert de
point d'appui discret~; le constat rejoint T-04 (identifiants par défaut).
La correction tient surtout au cloisonnement du réseau des objets.
\end{boxscenario}

\begin{boxobjet}{T-04 · Identifiants par défaut sur des équipements}
\textbf{Description~:} imprimantes et une borne gardent leur mot de passe
d'administration d'usine. \textbf{Preuve~:} accès aux interfaces
d'administration (sans modification). \textbf{Gravité~:} moyenne.
\textbf{Impact~:} prise de contrôle d'équipements, point d'appui sur le
réseau. \textbf{Recommandation~:} changer les mots de passe, cloisonner le
réseau des objets, inventorier. \textbf{Détection~:} connexions
d'administration depuis des adresses inhabituelles.
\end{boxobjet}

% ============================================================================
\section{Tester Légalement}
\label{sec:c18-legal}
% ============================================================================

\begin{boxprocedure}
\textbf{Encadrer un test sans-fil}
\begin{enumerate}
    \item Autorisation écrite précisant les réseaux visés et le lieu.
    \item Réseau et comptes dédiés au test pour tout essai de capture.
    \item Pas d'interception des communications de tiers.
    \item Inventaire des objets connectés avant les essais.
    \item Restitution des équipements dans leur état initial.
\end{enumerate}
\end{boxprocedure}

\begin{boxmetier}{Pentesteur et administrateur réseau}
\textbf{Cloisonner d'abord.} La meilleure réponse au risque sans-fil et
objets n'est pas l'audit répété, mais le cloisonnement~: un réseau séparé
pour les objets, un réseau invité coupé du réseau interne.
\end{boxmetier}

\separateur

\begin{boxresume}
\textbf{Ce qu'il faut retenir de ce chapitre~:}
\begin{itemize}
    \item \textbf{Le sans-fil franchit les murs}~; sa portée dépasse le
          bâtiment.
    \item \textbf{Point d'accès pirate} et \textbf{portail jumeau} sont les
          deux menaces majeures.
    \item \textbf{Les objets connectés} cumulent identifiants par défaut et
          absence de mise à jour.
    \item \textbf{Capter les communications d'autrui} est une infraction~;
          on teste sur des réseaux dédiés.
    \item \textbf{Le cloisonnement} est la première parade.
\end{itemize}
\end{boxresume}

\begin{boxinfo}
\textbf{Points de contrôle de l'élève.} Je sais~: citer les faiblesses du
sans-fil~; définir point d'accès pirate et portail jumeau~; décrire les
risques des objets connectés~; encadrer légalement un test sans-fil.
\end{boxinfo}

\bigskip

\begin{boxexo}[title={\ding{45}\quad Exercice~18.1. Portée \niveaubase}]
Pourquoi le réseau sans-fil de l'entrepôt est-il captable depuis le
parking~?
\end{boxexo}

\begin{boxexo}[title={\ding{45}\quad Exercice~18.2. Jumeau \niveaubase}]
Quelle est la différence entre un point d'accès pirate et un portail
jumeau~?
\end{boxexo}

\begin{boxexo}[title={\ding{45}\quad Exercice~18.3. Objets \niveaubase}]
Citez trois faiblesses récurrentes des objets connectés.
\end{boxexo}

\begin{boxexo}[title={\ding{45}\quad Exercice~18.4. Légalité \niveaumoyen}]
Pourquoi un test de portail jumeau sur les salariés réels est-il interdit~?
Comment le mener légalement~?
\end{boxexo}

\begin{boxexo}[title={\ding{45}\quad Exercice~18.5. T-04 \niveaumoyen}]
Rédigez la correction de T-04 pour l'entrepôt de Douala.
\end{boxexo}

\begin{boxexo}[title={\ding{45}\quad Exercice~18.6. Cloisonnement \niveaumoyen}]
Décrivez le cloisonnement réseau à mettre en place pour les objets
connectés.
\end{boxexo}

\begin{boxexo}[title={\ding{45}\quad Exercice~18.7. Réseau invité \niveaumoyen}]
Pourquoi un réseau invité relié au réseau interne est-il dangereux~?
\end{boxexo}

\begin{boxexo}[title={\ding{45}\quad Exercice~18.8. Inventaire \niveaumoyen}]
Pourquoi commence-t-on un test d'objets connectés par un inventaire~?
\end{boxexo}

\begin{boxexo}[title={\ding{45}\quad Exercice~18.9. Plan \niveauexpert}]
Rédigez le plan d'un test sans-fil autorisé pour l'entrepôt, de
l'autorisation à la restitution.
\end{boxexo}

\begin{boxexo}[title={\ding{45}\quad Exercice~18.10. Architecture \niveauexpert}]
Proposez une architecture réseau qui sépare bureautique, objets et invités.
\end{boxexo}

\bigskip

\begin{boxinfo}
\textbf{Pour aller plus loin.}
\begin{itemize}
    \item MITRE ATT\&CK~\cite{mitreattack}.
\end{itemize}
\end{boxinfo}

\begin{boxinfo}
\textbf{Transition.} La faiblesse la plus universelle n'est pas technique~:
le Chapitre~\ref{chap:social} traite de l'ingénierie sociale.
\end{boxinfo}
